\begin{figure}[t!]
    \centering    
    \includegraphics[width=0.48\textwidth]{figures/MAS4TS_comparison.pdf}
    \caption{The motivation of \model. Existing approaches (left) struggle to support intuitive reasoning and task-adaptive execution, while \model (right) organizes time series analysis into a unified paradigm that collaboratively analyzes, reasons, and executes.}
    \label{fig:motivation}
\vspace{-4mm}
\end{figure}

\section{Introduction}

Time series analysis is fundamental to a wide range of real-world applications, including energy systems~\cite{deb2017review}, transportation~\cite{zheng2020traffic}, and other critical domains~\cite{idrees2019prediction,karevan2020transductive,schneider1974climate}.
Core tasks such as forecasting, classification, imputation, and anomaly detection~\cite{lim2021time} play a crucial role in supporting data-driven decision-making in these scenarios.

Recent advances in pretrained Large Language Models (LLMs) have demonstrated strong capabilities in pattern understanding and general reasoning~\cite{brown2020language,touvron2023llama}.
Motivated by this progress, a growing body of work explores leveraging pretrained models for time series analysis~\cite{jin2023time,liu2024unitime,zhou2023one}, aiming to improve generalization and contextual modeling beyond traditional neural architectures~\cite{wu2023timesnet,zhou2021informer,nie2022time}.
In parallel, agentic methods and multi-agent systems (MAS)~\cite{ravuru2024agentic,gu2025argos,chang2025survey,wu2024autogen,hong2023metagpt,qian2023chatdev,huang2025many} are emerging as a practical paradigm for tool orchestration, task decomposition, and coordinated execution.

Despite these advances, the application of multi-agent paradigms to general time series analysis remains largely unexplored.
Most existing approaches primarily rely on implicit historical pattern matching in numerical or embedding spaces, lacking explicit and interpretable reasoning mechanisms that capture latent temporal structures.
Even recent vision-enhanced time series methods~\cite{chen2024visiontsvisualmaskedautoencoders,zhong2025time,ruan2025vision} often treat visual representations as auxiliary features, rather than leveraging the ability of modern Vision-Language Models (VLMs) to perform \emph{direct reasoning on time series plots} that naturally encode trends, regime shifts, and structural transitions.
Moreover, current methods are typically designed for isolated tasks under fixed inference pipelines, with limited support for adaptive tool orchestration and coordinated execution across heterogeneous time series workloads.
These limitations can be summarized by two conceptual gaps:
\begin{itemize}[leftmargin=*, itemsep=0pt]
\item \textbf{From Representation Learning to Structural Reasoning.}
Existing methods operate in numerical or latent embedding spaces where morphological patterns remain implicit and entangled.
Recent vision-enhanced methods (e.g., TimeVLM, VisionTS) convert time series into vision primarily for \emph{representation learning}, but fail to exploit visual reasoning capabilities to extract explicit structural knowledge from plots like human analysts naturally use for temporal understanding.
\textit{Can we shift from learning representations to performing intuitive structural reasoning over time series visualizations?}

\item \textbf{From Task-Specific Models to Tool-Driven Execution.}
Current paradigms treat forecasting, classification, imputation, and anomaly detection as separate problems, typically requiring task-specific heads and individual procedures even when sharing a common backbone.
Unlike multi-agent systems that leverage tool orchestration for adaptive problem-solving, time series methods operate under the \textit{fixed model then single-task finetuning} paradigms.
It prevents pretrained models from functioning as general-purpose, task-adaptive time series solvers and necessitates multiplicative scaling costs.
\textit{Can we unify diverse time series tasks through shared structural reasoning while enabling adaptive tool selection for task-specific execution?}
\end{itemize}

General time series tasks fundamentally require an understanding of morphological semantics that are explicit in visual plots yet implicit in numerical embeddings.
This reveals that the bottleneck of time series intelligence lies in structural reasoning rather than representation learning, suggesting a paradigm shift: unify tasks through visual reasoning while enabling adaptive tool-driven execution.

Based on this insight, we propose \textbf{MAS4TS}, the first tool-driven \underline{M}ulti-\underline{A}gent \underline{S}ystem for general \underline{T}ime \underline{S}eries analysis, built upon an \textbf{Analyzer--Reasoner--Executor} paradigm.
MAS4TS decomposes time series analysis into complementary roles: the \emph{Analyzer} performs statistical analysis and data preprocessing; the \emph{Reasoner} integrates visual and numerical reasoning to extract high-level temporal priors and derive latent trajectory constraints; and the \emph{Executor} invokes appropriate tools and models to generate task-specific outputs.
Specifically, MAS4TS performs visual reasoning directly on time series plots using VLMs to obtain structured anchors, which subsequently guide trajectory reconstruction in latent space.
Through shared memory and inter-agent communication, specialized agents collaborate to solve diverse time series tasks within a unified framework.
Our key contributions are summarized as follows:
\begin{itemize}[leftmargin=*, itemsep=0pt]
\item We propose \textbf{MAS4TS}, the first tool-driven multi-agent system for general time series analysis, establishing an \textit{Analyzer--Reasoner--Executor} paradigm for task-adaptive time series modeling.
\item We introduce an intuitive reasoning approach that utilizes VLMs to extract visual anchors from plots and priors then transforms them into numerical constraints for latent trajectory reconstruction.
\item We conduct extensive experiments and demonstrate the superior performance of \model among forecasting, classification, imputation, and anomaly detection tasks.
\end{itemize}
